\chapter{Abstimmen}

\section{Der Leitfall}

\begin{merksatz}
\boxtitel{Eine Gerätebeschaffung}
Das Klinikum kann eines von mehreren Großgeräten anschaffen. Fünf Abteilungen stimmen ab, jede reiht die Geräte nach ihrem Bedarf.

\begin{center}
\begin{minipage}{0.9\textwidth}
\textbf{Die Frage ist neu:} Nicht wie einzelne sich verhalten, sondern wie man aus vielen Reihungen eine gemeinsame macht.
\end{minipage}
\end{center}

\begin{center}
{\small
\renewcommand{\arraystretch}{1.5}
\begin{tabular}{@{}P{2.4cm}P{6.2cm}@{}}
\toprule
\textbf{Kapitel 11} & Man verändert Auszahlungen, um ein gewünschtes Verhalten herbeizuführen.\\
\textbf{Hier} & Die Auszahlungen liegen fest. Gesucht ist ein \kursiv{Verfahren}, das sie zusammenführt.\\
\bottomrule
\end{tabular}}
\end{center}

\begin{center}
\begin{minipage}{0.88\textwidth}
\kursiv{Man erwartet, dass sich ein vernünftiges Verfahren finden lässt. \textbf{Dieses Kapitel zeigt, dass keines existiert} --- und zwar nicht, weil man noch nicht genug gesucht hätte.}
\end{minipage}
\end{center}
\end{merksatz}

\section{Das Paradoxon von Condorcet}

\begin{merksatz}
\boxtitel{Drei Abteilungen, drei Geräte}
\begin{center}
{\small
\renewcommand{\arraystretch}{1.5}
\begin{tabular}{@{}P{2.4cm}P{4.0cm}@{}}
\toprule
\textbf{Abteilung} & \textbf{Reihung}\\
\midrule
Chirurgie & $A \succ B \succ C$\\
Innere & $B \succ C \succ A$\\
Radiologie & $C \succ A \succ B$\\
\bottomrule
\end{tabular}}
\end{center}

\schritt{Die Paarvergleiche}
\begin{center}
{\small
\renewcommand{\arraystretch}{1.5}
\begin{tabular}{@{}P{2.2cm}P{2.0cm}P{2.6cm}@{}}
\toprule
\textbf{Vergleich} & \textbf{Ergebnis} & \\
\midrule
$A$ gegen $B$ & $2:1$ & $A$ gewinnt\\
$B$ gegen $C$ & $2:1$ & $B$ gewinnt\\
$C$ gegen $A$ & $2:1$ & $C$ gewinnt\\
\bottomrule
\end{tabular}}
\end{center}

\begin{center}
\begin{tikzpicture}[scale=1.0]
\node[knotenb] (A) at (0,1.3) {$A$};
\node[knotenb] (B) at (1.5,-0.8) {$B$};
\node[knotenb] (C) at (-1.5,-0.8) {$C$};
\draw[kanted, -{Stealth[length=4.5pt]}] (A)--(B);
\draw[kanted, -{Stealth[length=4.5pt]}] (B)--(C);
\draw[kanted, -{Stealth[length=4.5pt]}] (C)--(A);
\node[font=\scriptsize\itshape, color=grau] at (0,-1.6) {jede Mehrheit ist klar, die Reihung ist es nicht};
\end{tikzpicture}
\end{center}

\begin{center}
\begin{minipage}{0.9\textwidth}
\textbf{Die Mehrheitsreihung ist nicht übertragbar.} Es gibt keinen Sieger, obwohl jeder Einzelvergleich eindeutig ausgeht.
\end{minipage}
\end{center}

\begin{center}
\begin{minipage}{0.88\textwidth}
\kursiv{Die Übertragbarkeit war in Satz 4.1 eine Forderung an die Vorlieben des Einzelnen --- und jede Abteilung erfüllt sie hier. \textbf{Die Mehrheitsregel erhält sie nicht.} Das ist der Kern des Kapitels: Eine Eigenschaft, die alle Beteiligten haben, muss das Ergebnis der Zusammenführung nicht haben.}
\end{minipage}
\end{center}

\schritt{Die Folge für die Tagesordnung}
Stimmt man paarweise ab, so entscheidet die Reihenfolge der Abstimmungen.

\begin{center}
{\small
\renewcommand{\arraystretch}{1.5}
\begin{tabular}{@{}P{3.0cm}P{2.6cm}P{2.4cm}@{}}
\toprule
\textbf{Tagesordnung} & \textbf{erster Sieger} & \textbf{Ergebnis}\\
\midrule
erst $A$/$B$, dann gegen $C$ & $A$ & $C$\\
erst $B$/$C$, dann gegen $A$ & $B$ & $A$\\
erst $C$/$A$, dann gegen $B$ & $C$ & $B$\\
\bottomrule
\end{tabular}}
\end{center}

\begin{center}
\begin{minipage}{0.9\textwidth}
\textbf{Wer die Tagesordnung bestimmt, bestimmt das Ergebnis --- vollständig.}
\end{minipage}
\end{center}

\begin{center}
\begin{minipage}{0.88\textwidth}
\kursiv{Und zwar ohne jede Regelverletzung: Jede einzelne Abstimmung geht mit klarer Mehrheit aus. \textbf{Die Macht liegt bei der Sitzungsleitung, nicht bei den Stimmen} --- ein Befund, den jeder kennt, der Gremienarbeit betreibt, und der hier eine exakte Begründung erhält.}
\end{minipage}
\end{center}
\end{merksatz}

\section{Vier Verfahren, vier Sieger}

\begin{merksatz}
\boxtitel{Ein Beispiel mit fünf Geräten}
Dreiundfünfzig Stimmberechtigte, fünf Geräte:

\begin{center}
{\small
\renewcommand{\arraystretch}{1.5}
\begin{tabular}{@{}P{1.6cm}P{5.0cm}@{}}
\toprule
\textbf{Anzahl} & \textbf{Reihung}\\
\midrule
$18$ & $A \succ D \succ E \succ C \succ B$\\
$12$ & $B \succ E \succ D \succ C \succ A$\\
$10$ & $C \succ B \succ E \succ D \succ A$\\
$9$ & $D \succ C \succ E \succ B \succ A$\\
$4$ & $E \succ B \succ D \succ C \succ A$\\
\bottomrule
\end{tabular}}
\end{center}

\schritt{Verfahren 1 --- relative Mehrheit}
Nur der erste Platz zählt.
\begin{center}
{\small
\renewcommand{\arraystretch}{1.44}
\begin{tabular}{@{}P{1.2cm}P{1.2cm}P{1.2cm}P{1.2cm}P{1.2cm}@{}}
\toprule
$A$ & $B$ & $C$ & $D$ & $E$\\
$\mathbf{18}$ & $12$ & $10$ & $9$ & $4$\\
\bottomrule
\end{tabular}}
\end{center}
\[ \Longrightarrow\ A \]

\schritt{Verfahren 2 --- Stichwahl der beiden Besten}
$A$ gegen $B$:
\[ A: 18, \qquad B: 12+10+9+4 = 35 \]
\[ \Longrightarrow\ B \]

\schritt{Verfahren 3 --- schrittweise Ausscheidung}
Immer der Letzte fällt heraus, seine Stimmen gehen weiter.
\begin{center}
{\small
\renewcommand{\arraystretch}{1.44}
\begin{tabular}{@{}P{1.6cm}P{1.2cm}P{1.2cm}P{1.2cm}P{1.2cm}P{1.2cm}@{}}
\toprule
\textbf{Runde} & $A$ & $B$ & $C$ & $D$ & $E$\\
\midrule
$1$ & $18$ & $12$ & $10$ & $9$ & $4$\\
$2$ & $18$ & $16$ & $10$ & $9$ & ---\\
$3$ & $18$ & $16$ & $19$ & --- & ---\\
$4$ & $18$ & --- & $\mathbf{35}$ & & \\
\bottomrule
\end{tabular}}
\end{center}
\[ \Longrightarrow\ C \]

\schritt{Verfahren 4 --- Borda-Punkte}
Vier Punkte für den ersten Platz, drei für den zweiten, und so fort.
\begin{center}
{\small
\renewcommand{\arraystretch}{1.44}
\begin{tabular}{@{}P{1.2cm}P{1.4cm}P{1.4cm}P{1.4cm}P{1.4cm}P{1.4cm}@{}}
\toprule
& $A$ & $B$ & $C$ & $D$ & $E$\\
\midrule
Punkte & $72$ & $99$ & $101$ & $\mathbf{132}$ & $126$\\
\bottomrule
\end{tabular}}
\end{center}

\kursiv{Zur Spalte $D$:} $18\cdot3 + 12\cdot2 + 10\cdot1 + 9\cdot4 + 4\cdot2 = 54+24+10+36+8 = 132$

\[ \Longrightarrow\ D \]

\schritt{Verfahren 5 --- paarweise Mehrheiten}
$D$ schlägt jeden anderen:
\begin{center}
{\small
\renewcommand{\arraystretch}{1.46}
\begin{tabular}{@{}P{2.0cm}P{2.0cm}P{2.6cm}@{}}
\toprule
\textbf{Vergleich} & \textbf{Ergebnis} & \\
\midrule
$D$ gegen $A$ & $35:18$ & \\
$D$ gegen $B$ & $27:26$ & knapp\\
$D$ gegen $C$ & $43:10$ & \\
$D$ gegen $E$ & $27:26$ & knapp\\
\bottomrule
\end{tabular}}
\end{center}
\[ \Longrightarrow\ D \ \text{(Condorcet-Sieger)} \]

\begin{center}
\begin{minipage}{0.9\textwidth}
\textbf{Vier verschiedene Sieger bei fünf gebräuchlichen Verfahren --- aus denselben Stimmzetteln.}
\end{minipage}
\end{center}

\begin{center}
\begin{minipage}{0.88\textwidth}
\kursiv{Besonders unangenehm ist $A$: Es gewinnt die relative Mehrheit und verliert jeden Paarvergleich mit $18$ gegen $35$. \textbf{Achtzehn von dreiundfünfzig wollen $A$; fünfunddreißig halten es für das schlechteste.} Die relative Mehrheitswahl ist das einzige Verfahren, das diesen Umstand vollständig ignoriert --- und das am weitesten verbreitete.}
\end{minipage}
\end{center}
\end{merksatz}

\section{Der Satz von Arrow}

\begin{merksatz}
\boxtitel{Was man von einem Verfahren erwartet}
Ein \Begriff{Wahlverfahren} bildet die Reihungen aller Beteiligten auf eine gemeinsame Reihung ab.

\begin{center}
{\small
\renewcommand{\arraystretch}{1.55}
\begin{tabular}{@{}P{2.6cm}P{6.0cm}@{}}
\toprule
\textbf{Universalität} & Jede Zusammenstellung von Reihungen ist zulässig, und das Ergebnis ist stets eine vollständige, übertragbare Reihung.\\
\textbf{Einstimmigkeit} & Ziehen alle $x$ dem $y$ vor, so auch das Ergebnis.\\
\textbf{Unabhängigkeit} & Ob $x$ vor $y$ steht, hängt nur davon ab, wie die Beteiligten $x$ und $y$ zueinander reihen --- nicht von einer dritten Möglichkeit $z$.\\
\textbf{Keine Diktatur} & Es gibt keinen Beteiligten, dessen Reihung stets das Ergebnis ist.\\
\bottomrule
\end{tabular}}
\end{center}
\end{merksatz}

\begin{merksatz}
\boxtitel{Satz 12.1 --- Arrow, 1951}
Bei mindestens drei Möglichkeiten und mindestens zwei Beteiligten gibt es kein Verfahren, das alle vier Bedingungen erfüllt.

\begin{center}
\begin{minipage}{0.9\textwidth}
\textbf{Nicht: man kennt keines. Es gibt keines.}
\end{minipage}
\end{center}

\begin{center}
\begin{minipage}{0.88\textwidth}
\kursiv{Der Beweis ist nicht lang, aber trickreich; er führt über den Begriff einer entscheidenden Gruppe und zeigt, dass diese sich auf eine einzige Person zusammenziehen lässt. \textbf{Wir verzichten auf ihn} --- die zweite der neun Auslassungen dieses Bandes. Die Fundstelle ist Osborne und Rubinstein\zit{19}, Kapitel~10.}
\end{minipage}
\end{center}

\schritt{Wie die Verfahren scheitern}
\begin{center}
{\small
\renewcommand{\arraystretch}{1.5}
\begin{tabular}{@{}P{2.4cm}P{6.2cm}@{}}
\toprule
\textbf{Mehrheitsregel} & Verletzt die Universalität --- das Ergebnis ist nicht immer übertragbar (Abschnitt 12.2).\\
\textbf{Borda} & Verletzt die Unabhängigkeit --- eine zusätzliche Möglichkeit verschiebt die Punkte.\\
\textbf{Diktatur} & Erfüllt die ersten drei und ist ausgeschlossen.\\
\bottomrule
\end{tabular}}
\end{center}

\schritt{Die Unabhängigkeit ist die strittige Bedingung}
\begin{center}
\begin{minipage}{0.88\textwidth}
\kursiv{Sie verlangt, dass die Stärke einer Vorliebe unberücksichtigt bleibt --- ob jemand $x$ knapp oder deutlich vor $y$ setzt, darf keine Rolle spielen. \textbf{Genau das macht Borda anders}, und genau deshalb scheitert es an dieser Bedingung. Wer sie aufgibt, muss allerdings Nutzen zwischen Personen vergleichen, und Satz 4.1 sagt, dass das keinen Sinn ergibt.}
\end{minipage}
\end{center}
\end{merksatz}

\begin{merksatz}
\boxtitel{Satz 12.2 --- Gibbard und Satterthwaite, 1973/75}
Jedes Verfahren, das bei mindestens drei Möglichkeiten jede von ihnen erreichen kann und keine Diktatur ist, lässt sich durch falsche Angaben zum eigenen Vorteil beeinflussen.

\begin{center}
\begin{minipage}{0.9\textwidth}
\textbf{Es gibt kein Verfahren, bei dem ehrliches Abstimmen stets die beste Wahl ist.}
\end{minipage}
\end{center}

\begin{center}
\begin{minipage}{0.88\textwidth}
\kursiv{Der Satz ist eine Folgerung aus Satz 12.1 und in gewisser Weise die praktisch schärfere Aussage. \textbf{Arrow sagt, dass die Zusammenführung Mängel hat; Gibbard und Satterthwaite sagen, dass die Stimmzettel selbst unzuverlässig sind.} Wer taktisch abstimmt, gibt nicht seine Vorlieben an, und damit fehlt dem Verfahren die Eingabe, die es verarbeiten soll.}
\end{minipage}
\end{center}

\schritt{Borda wusste es}
\begin{center}
\begin{minipage}{0.88\textwidth}
\kursiv{Auf den Einwand, sein Verfahren lasse sich manipulieren, soll \Person{Jean-Charles de Borda} geantwortet haben, es sei nur für ehrliche Menschen gedacht. \textbf{Das ist keine Verteidigung, sondern eine Beschreibung der Lage} --- und Satz 12.2 besagt, dass es keinem anderen Verfahren besser ergeht.}
\end{minipage}
\end{center}
\end{merksatz}

\section{Der Bruchfall}

\begin{warnung}
\boxtitelr{Was die Sätze nicht bedeuten}
\begin{center}
{\small
\renewcommand{\arraystretch}{1.52}
\begin{tabular}{@{}P{2.4cm}P{6.2cm}@{}}
\toprule
\textbf{Nicht:} & \enquote{Alle Verfahren sind gleich schlecht.} Sie scheitern an verschiedenen Bedingungen und in verschieden häufigen Fällen.\\
\textbf{Nicht:} & \enquote{Abstimmen ist sinnlos.} Bei zwei Möglichkeiten gilt keiner der Sätze --- die einfache Mehrheit erfüllt alles.\\
\textbf{Nicht:} & \enquote{Manipulation ist häufig.} Sie ist möglich; ob sie gelingt, hängt davon ab, wie viel man über die anderen weiß.\\
\bottomrule
\end{tabular}}
\end{center}

\schritt{Zur zweiten Zeile}
\Person{Kenneth May} zeigte 1952, dass bei genau zwei Möglichkeiten die einfache Mehrheit das einzige Verfahren ist, das Gleichbehandlung der Stimmen, Gleichbehandlung der Möglichkeiten und Monotonie erfüllt.

\begin{center}
\begin{minipage}{0.88\textwidth}
\kursiv{\textbf{Die Schwierigkeit beginnt also exakt bei drei.} Das erklärt zugleich, warum Stichwahlen und Ausschussverfahren so verbreitet sind: Sie zerlegen die Aufgabe in Entscheidungen zwischen zweien --- um den Preis, dass die Reihenfolge das Ergebnis bestimmt, wie Abschnitt 12.2 zeigt.}
\end{minipage}
\end{center}

\schritt{Wie oft es überhaupt schiefgeht}
\begin{center}
{\small
\renewcommand{\arraystretch}{1.48}
\begin{tabular}{@{}P{2.6cm}P{2.2cm}P{3.0cm}@{}}
\toprule
\textbf{Möglich\-keiten} & \textbf{kein Sieger} & \\
\midrule
$3$ & $\approx 9\,\%$ & bei zufälligen Reihungen\\
$5$ & $\approx 25\,\%$ & \\
$10$ & $\approx 50\,\%$ & \\
\bottomrule
\end{tabular}}
\end{center}

\begin{center}
\begin{minipage}{0.88\textwidth}
\kursiv{Die Werte gelten für rein zufällige, voneinander unabhängige Reihungen und viele Stimmen --- der ungünstigste Fall. \textbf{In tatsächlichen Wahlen sind die Reihungen nicht zufällig}, sondern häufig entlang einer gemeinsamen Achse geordnet, und dann gibt es fast immer einen Sieger. Abschnitt 12.6 behandelt diesen Ausweg.}
\end{minipage}
\end{center}
\end{warnung}

\begin{merksatz}
\boxtitel{Der Ausweg: eingipflige Vorlieben}
Lassen sich die Möglichkeiten auf einer Linie anordnen, so dass jeder Beteiligte einen Bestpunkt hat und nach beiden Seiten hin abfällt, so heißen die Vorlieben \Begriff{eingipflig}.

\schritt{Satz 12.3 --- der Medianwähler, Black 1948}
Bei eingipfligen Vorlieben und ungerader Anzahl Beteiligter ist der Bestpunkt des mittleren Beteiligten ein Condorcet-Sieger.

\schritt{Beweis}
Sei $m$ der Median der Bestpunkte. Für jede Möglichkeit $y$ rechts von $m$ ziehen alle, deren Bestpunkt bei $m$ oder links davon liegt, $m$ dem $y$ vor --- das sind mindestens die Hälfte plus einer.

Ebenso für $y$ links von $m$. $\blacksquare$

\begin{center}
\begin{minipage}{0.9\textwidth}
\textbf{Eine Anordnung auf einer Linie genügt, und die Schwierigkeit verschwindet.}
\end{minipage}
\end{center}

\begin{center}
\begin{minipage}{0.88\textwidth}
\kursiv{Das erklärt, warum Abstimmungen über Beträge --- Budgethöhe, Personalschlüssel, Schwellenwerte --- meist glatt verlaufen: Sie sind von Natur aus eingipflig. \textbf{Schwierig wird es, wenn die Möglichkeiten sich in mehreren Hinsichten unterscheiden} --- etwa bei der Gerätewahl, wo Anschaffungspreis, Platzbedarf und Fachgebiet gegeneinanderstehen. Bei zwei oder mehr Dimensionen gilt der Satz nicht mehr.}
\end{minipage}
\end{center}
\end{merksatz}

\section{Übungsaufgaben}

\begin{aufgabe}
\boxtitel{Aufgabe 12.1 --- Ein Verfahren wählen}
Prüfen Sie am Beispiel aus Abschnitt 12.3, welches Gerät jedes Verfahren wählt, und beurteilen Sie die Ergebnisse.

\loesung
\schritt{Schritt 1 --- Die Übersicht}
\begin{center}
{\small
\renewcommand{\arraystretch}{1.5}
\begin{tabular}{@{}P{3.0cm}P{1.4cm}P{3.0cm}@{}}
\toprule
\textbf{Verfahren} & \textbf{Sieger} & \\
\midrule
relative Mehrheit & $A$ & \\
Zweierstichwahl & $B$ & \\
Ausscheidung & $C$ & \\
Borda & $D$ & \\
Paarvergleich & $D$ & Condorcet-Sieger\\
\bottomrule
\end{tabular}}
\end{center}

\schritt{Schritt 2 --- $A$ beurteilen}
\[ A \ \text{gegen jeden anderen}: 18:35 \]

\begin{center}
\begin{minipage}{0.9\textwidth}
\textbf{$A$ verliert jeden Paarvergleich --- es ist ein Condorcet-Verlierer.}
\end{minipage}
\end{center}

\kursiv{Grund:} Alle $35$ Stimmberechtigten, die $A$ nicht auf Platz eins setzen, setzen es auf Platz fünf.

\schritt{Schritt 3 --- $C$ beurteilen}
\begin{center}
\renewcommand{\arraystretch}{1.48}
\begin{tabular}{@{}P{2.4cm}P{2.0cm}P{2.6cm}@{}}
\toprule
\textbf{Vergleich} & \textbf{Ergebnis} & \\
\midrule
$C$ gegen $D$ & $10:43$ & deutlich verloren\\
$C$ gegen $E$ & $19:34$ & verloren\\
$C$ gegen $B$ & $37:16$ & gewonnen\\
\bottomrule
\end{tabular}
\end{center}

\begin{center}
\begin{minipage}{0.9\textwidth}
\textbf{Das Ausscheidungsverfahren wählt ein Gerät, das gegen $D$ mit $10$ zu $43$ verliert.}
\end{minipage}
\end{center}

\begin{center}
\begin{minipage}{0.88\textwidth}
\kursiv{Das ist kein Randfall dieses Beispiels, sondern eine bekannte Eigenschaft: \textbf{Das Ausscheidungsverfahren kann den Condorcet-Sieger in einer frühen Runde eliminieren}, weil es nur erste Plätze zählt. $D$ hatte nur neun erste Plätze und wurde in Runde 3 ausgeschieden --- obwohl es gegen alle anderen eine Mehrheit hat.}
\end{minipage}
\end{center}

\schritt{Schritt 4 --- $D$ beurteilen}
\begin{center}
\begin{minipage}{0.9\textwidth}
\textbf{$D$ schlägt jeden anderen in der direkten Abstimmung --- zweimal allerdings nur mit $27$ zu $26$.}
\end{minipage}
\end{center}

\begin{center}
\begin{minipage}{0.88\textwidth}
\kursiv{Die Knappheit gegen $B$ und gegen $E$ ist bemerkenswert. \textbf{Ein Condorcet-Sieger ist nicht notwendig ein starker Sieger}, sondern nur einer, der keinen Paarvergleich verliert. Eine einzige Stimmänderung könnte die Lage kippen --- und dann gäbe es keinen Condorcet-Sieger mehr.}
\end{minipage}
\end{center}

\schritt{Schritt 5 --- Die Empfehlung}
\begin{center}
{\small
\renewcommand{\arraystretch}{1.5}
\begin{tabular}{@{}P{2.4cm}P{6.2cm}@{}}
\toprule
\textbf{Wenn es einen gibt} & Der Condorcet-Sieger ist die sicherste Wahl --- keine Mehrheit zieht etwas anderes vor.\\
\textbf{Wenn nicht} & Borda liefert ein plausibles Ergebnis und schließt den Condorcet-Verlierer aus; manipulierbar bleibt es.\\
\textbf{Nie} & Die relative Mehrheit bei mehr als zwei Möglichkeiten.\\
\bottomrule
\end{tabular}}
\end{center}

\begin{center}
\begin{minipage}{0.88\textwidth}
\kursiv{Der letzte Punkt ist die einzige klare Handlungsempfehlung, die dieses Kapitel hergibt. \textbf{Sie folgt nicht aus Arrow, sondern aus der Anschauung}: Ein Verfahren, das den Condorcet-Verlierer zum Sieger erklären kann, ist schwer zu rechtfertigen.}
\end{minipage}
\end{center}
\end{aufgabe}

\begin{aufgabe}
\boxtitel{Aufgabe 12.2 --- Die Tagesordnung nutzen}
Im Paradoxon aus Abschnitt 12.2 leitet die Radiologie die Sitzung und möchte, dass $C$ gewinnt. Welche Tagesordnung wählt sie?

\loesung
\schritt{Schritt 1 --- Die Zyklusrichtung}
\[ A \succ B, \qquad B \succ C, \qquad C \succ A \]

\schritt{Schritt 2 --- Der Gedanke}
Die eigene Möglichkeit soll \kursiv{zuletzt} antreten, gegen denjenigen, den sie schlägt.

\[ C \ \text{schlägt} \ A \quad\Longrightarrow\quad A \ \text{muss die erste Runde gewinnen} \]

\schritt{Schritt 3 --- Die Tagesordnung}
\begin{center}
{\small
\renewcommand{\arraystretch}{1.5}
\begin{tabular}{@{}P{1.8cm}P{2.8cm}P{2.6cm}@{}}
\toprule
\textbf{Runde} & \textbf{Abstimmung} & \textbf{Sieger}\\
\midrule
$1$ & $A$ gegen $B$ & $A$ $(2:1)$\\
$2$ & $A$ gegen $C$ & $\mathbf{C}$ $(2:1)$\\
\bottomrule
\end{tabular}}
\end{center}

\schritt{Schritt 4 --- Die Probe für die anderen}
\begin{center}
{\small
\renewcommand{\arraystretch}{1.5}
\begin{tabular}{@{}P{2.6cm}P{2.4cm}P{2.4cm}@{}}
\toprule
\textbf{Wer leitet} & \textbf{Erste Runde} & \textbf{Ergebnis}\\
\midrule
Chirurgie ($A$) & $B$ gegen $C$ & $A$\\
Innere ($B$) & $C$ gegen $A$ & $B$\\
Radiologie ($C$) & $A$ gegen $B$ & $C$\\
\bottomrule
\end{tabular}}
\end{center}

\begin{center}
\begin{minipage}{0.9\textwidth}
\textbf{Jede Abteilung kann ihr eigenes Gerät durchsetzen, wenn sie die Tagesordnung bestimmt.}
\end{minipage}
\end{center}

\schritt{Schritt 5 --- Die Regel}
\begin{center}
\begin{minipage}{0.9\textwidth}
\textbf{Die eigene Möglichkeit in die letzte Runde --- und als Gegner denjenigen, den sie schlägt.}
\end{minipage}
\end{center}

\begin{center}
\begin{minipage}{0.88\textwidth}
\kursiv{Wer zuerst antritt, muss zweimal gewinnen; wer zuletzt antritt, nur einmal. \textbf{Das ist der Spätzieher-Vorteil} --- und damit das Gegenstück zum Erstzieher-Vorteil aus Kapitel~7. Ob Früh oder Spät nützt, hängt daran, ob man festlegt oder antwortet.}
\end{minipage}
\end{center}

\schritt{Schritt 6 --- Die Gegenmaßnahme}
\begin{center}
{\small
\renewcommand{\arraystretch}{1.5}
\begin{tabular}{@{}P{2.4cm}P{6.2cm}@{}}
\toprule
\textbf{Alle Paare} & Man stimmt über sämtliche Paare ab und sucht den Condorcet-Sieger. Gibt es keinen, ist das sichtbar.\\
\textbf{Borda} & Kennt keine Tagesordnung.\\
\textbf{Losentscheid} & Bei nachgewiesenem Zyklus --- unbefriedigend, aber unparteiisch.\\
\bottomrule
\end{tabular}}
\end{center}

\begin{center}
\begin{minipage}{0.88\textwidth}
\kursiv{Die erste Zeile ist bei wenigen Möglichkeiten praktikabel: Bei $k$ Möglichkeiten sind $\binom k2$ Abstimmungen nötig, bei fünf also zehn. \textbf{Der Aufwand verhindert die Manipulation der Reihenfolge vollständig} --- und macht zugleich sichtbar, wenn kein Sieger existiert. Das ist mehr wert als ein Ergebnis, das nur so aussieht, als wäre es eines.}
\end{minipage}
\end{center}
\end{aufgabe}

\begin{aufgabe}
\boxtitel{Aufgabe 12.3 --- Taktisch abstimmen}
Beim Borda-Verfahren im kleinen Beispiel --- $5$ mit $A\succ B\succ C$, $4$ mit $B\succ C\succ A$, $3$ mit $C\succ B\succ A$ --- können die $A$-Anhänger ihre Stimmzettel ändern. Was erreichen sie?

\loesung
\schritt{Schritt 1 --- Das ehrliche Ergebnis}
Punkte $2$, $1$, $0$:
\begin{center}
{\small
\renewcommand{\arraystretch}{1.48}
\begin{tabular}{@{}P{1.6cm}P{3.2cm}P{1.6cm}@{}}
\toprule
& \textbf{Rechnung} & \textbf{Summe}\\
\midrule
$A$ & $5\cdot2$ & $10$\\
$B$ & $5\cdot1+4\cdot2+3\cdot1$ & $\mathbf{16}$\\
$C$ & $4\cdot1+3\cdot2$ & $10$\\
\bottomrule
\end{tabular}}
\end{center}

$B$ gewinnt; die $A$-Anhänger erhalten ihre zweite Wahl.

\schritt{Schritt 2 --- Der Eingriff}
Die fünf $A$-Anhänger geben statt $A\succ B\succ C$ nun $A\succ C\succ B$ an --- sie stufen den gefährlichsten Gegner herab.

\begin{center}
{\small
\renewcommand{\arraystretch}{1.48}
\begin{tabular}{@{}P{1.6cm}P{3.2cm}P{1.6cm}@{}}
\toprule
& \textbf{Rechnung} & \textbf{Summe}\\
\midrule
$A$ & $5\cdot2$ & $10$\\
$B$ & $5\cdot0+4\cdot2+3\cdot1$ & $11$\\
$C$ & $5\cdot1+4\cdot1+3\cdot2$ & $\mathbf{15}$\\
\bottomrule
\end{tabular}}
\end{center}

\begin{center}
\begin{minipage}{0.9\textwidth}
\textbf{Nun gewinnt $C$ --- die \kursiv{letzte} Wahl der $A$-Anhänger.} Der Eingriff hat ihnen geschadet.
\end{minipage}
\end{center}

\schritt{Schritt 3 --- Der bessere Eingriff}
Um $A$ zum Sieg zu verhelfen, müssten die $A$-Anhänger $A$ um mindestens sieben Punkte verbessern --- das geht nicht, denn $A$ steht bei ihnen bereits ganz oben.

\begin{center}
\begin{minipage}{0.9\textwidth}
\textbf{$A$ ist nicht durchsetzbar.} Die einzige Frage ist, ob $B$ oder $C$ gewinnt.
\end{minipage}
\end{center}

\kursiv{Da die $A$-Anhänger $B$ dem $C$ vorziehen, ist ehrliches Abstimmen hier die beste Wahl.}

\schritt{Schritt 4 --- Wo der Eingriff lohnt}
Die $C$-Anhänger ($3$ Stimmen) geben statt $C\succ B\succ A$ nun $C\succ A\succ B$ an:
\begin{center}
{\small
\renewcommand{\arraystretch}{1.48}
\begin{tabular}{@{}P{1.6cm}P{3.2cm}P{1.6cm}@{}}
\toprule
& \textbf{Rechnung} & \textbf{Summe}\\
\midrule
$A$ & $5\cdot2+3\cdot1$ & $13$\\
$B$ & $5\cdot1+4\cdot2$ & $13$\\
$C$ & $4\cdot1+3\cdot2$ & $10$\\
\bottomrule
\end{tabular}}
\end{center}

\begin{center}
\begin{minipage}{0.9\textwidth}
\textbf{Gleichstand zwischen $A$ und $B$ --- und die $C$-Anhänger ziehen $B$ dem $A$ vor.} Auch dieser Eingriff schadet.
\end{minipage}
\end{center}

\schritt{Schritt 5 --- Der Fall, in dem es gelingt}
Die vier $B$-Anhänger stufen $C$ herab: statt $B\succ C\succ A$ nun $B\succ A\succ C$.
\begin{center}
{\small
\renewcommand{\arraystretch}{1.48}
\begin{tabular}{@{}P{1.6cm}P{3.2cm}P{1.6cm}@{}}
\toprule
& \textbf{Rechnung} & \textbf{Summe}\\
\midrule
$A$ & $5\cdot2+4\cdot1$ & $14$\\
$B$ & $5\cdot1+4\cdot2+3\cdot1$ & $\mathbf{16}$\\
$C$ & $3\cdot2$ & $6$\\
\bottomrule
\end{tabular}}
\end{center}

$B$ gewinnt weiterhin --- hier war der Eingriff unnötig, aber ungefährlich.

\schritt{Schritt 6 --- Die Lehre}
\begin{center}
\begin{minipage}{0.88\textwidth}
\kursiv{Satz 12.2 besagt, dass es \kursiv{Stimmenverteilungen} gibt, bei denen sich falsche Angaben lohnen --- nicht, dass es bei jeder Verteilung so ist. \textbf{In diesem Beispiel schadet der Eingriff jedem, der ihn versucht}, weil $B$ eine so breite Zustimmung hat. Manipulierbarkeit ist eine Eigenschaft des Verfahrens; ob sich Manipulation lohnt, hängt an den Stimmen --- und an dem, was man über sie weiß.}
\end{minipage}
\end{center}
\end{aufgabe}

\begin{geschichte}
\boxtitelg{Zwei Franzosen und ein Amerikaner}
\Person{Jean-Charles de Borda} (1733--1799) schlug sein Punkteverfahren 1770 der Académie des Sciences vor, die es zur Wahl ihrer Mitglieder übernahm.

\begin{center}
\begin{minipage}{0.9\textwidth}
\kursiv{Die Akademie behielt es bis 1800 bei. Dann wurde es auf Betreiben eines neuen Mitglieds abgeschafft --- \Person{Napoleon Bonaparte}, der als Erster Konsul den Vorsitz übernommen hatte. \textbf{Die Begründung ist nicht überliefert.}}
\end{minipage}
\end{center}

\schritt{Condorcet}
\Person{Nicolas de Condorcet} (1743--1794) wies 1785 auf das nach ihm benannte Paradoxon hin und schlug die paarweise Abstimmung vor.

\begin{center}
\begin{minipage}{0.9\textwidth}
\kursiv{Er war Mathematiker, Philosoph und Politiker, verfasste einen Verfassungsentwurf für die Republik und setzte sich früh für die Gleichberechtigung von Frauen und für die Abschaffung der Sklaverei ein. \textbf{Er starb 1794 in einer Gefängniszelle}, nachdem er als Girondist geächtet worden war --- die Umstände sind ungeklärt.}
\end{minipage}
\end{center}

\schritt{Die Wiederentdeckung}
Beide Arbeiten gerieten in Vergessenheit. \Person{Charles Dodgson} --- besser bekannt als \Person{Lewis Carroll} --- entwickelte in den 1870er Jahren unabhängig verwandte Verfahren für die Wahlen seines Colleges.

\begin{center}
\begin{minipage}{0.9\textwidth}
\kursiv{\Person{Duncan Black} stellte in den 1940er Jahren die Zusammenhänge her und bewies den Satz über den Medianwähler. \textbf{Er entdeckte dabei Condorcets und Bordas Arbeiten wieder}, die über anderthalb Jahrhunderte unbeachtet geblieben waren.}
\end{minipage}
\end{center}

\schritt{Arrow}
\Person{Kenneth Arrow} (1921--2017) bewies seinen Satz 1950 als Doktorand, im Alter von neunundzwanzig Jahren.

\begin{center}
\begin{minipage}{0.9\textwidth}
\kursiv{Er hatte nach einem Verfahren gesucht und stattdessen bewiesen, dass es keines gibt. \textbf{Die Arbeit umfasst in Buchform gut hundert Seiten} und begründete ein eigenes Fach, die Theorie der gesellschaftlichen Wahl. 1972 erhielt Arrow den Wirtschaftsnobelpreis --- als bis heute jüngster Träger.}
\end{minipage}
\end{center}
\end{geschichte}

\begin{forschung}
\boxtitelp{Wahlverfahren heute}
\begin{center}
{\small
\renewcommand{\arraystretch}{1.45}
\begin{tabular}{@{}P{3.2cm}P{5.0cm}@{}}
\toprule
\textbf{Im Einsatz} & Ausscheidungsverfahren in Australien und Irland, seit 2018 auch in mehreren amerikanischen Städten.\\
\textbf{Zustimmungs\-wahl} & Jeder kreuzt beliebig viele an; umgeht Arrow, weil keine Reihung verlangt wird.\\
\textbf{Rechenaufwand} & Manche Verfahren sind nachweislich schwer manipulierbar --- ein Schutz durch Komplexität.\\
\bottomrule
\end{tabular}}
\end{center}

\schritt{Zur zweiten Zeile}
Bei der Zustimmungswahl gibt jeder an, welche Möglichkeiten er für annehmbar hält.

\begin{center}
\begin{minipage}{0.88\textwidth}
\kursiv{Satz 12.1 setzt voraus, dass Reihungen die Eingabe sind --- daher greift er hier nicht unmittelbar. \textbf{Dafür entsteht ein neues Problem: Wo setzt man die Grenze?} Die Antwort hängt davon ab, was man über die anderen erwartet, und damit ist das taktische Element zurück. Satz 12.2 gilt für die Zustimmungswahl ebenfalls.}
\end{minipage}
\end{center}

\schritt{Zur dritten Zeile}
Für einige Verfahren ist die Berechnung einer wirksamen Manipulation $\mathsf{NP}$-schwer.

\begin{center}
\begin{minipage}{0.88\textwidth}
\kursiv{Das klingt nach einem Ausweg und ist in der Praxis schwächer, als es aussieht: \textbf{Die Schwere gilt für den schlechtesten Fall}, und die häufig auftretenden Stimmenverteilungen sind oft leicht zu manipulieren. Es ist dasselbe Problem wie in Band 20, Kapitel~10 --- ein Verfahren, dessen Sicherheit an einer Worst-Case-Aussage hängt, ist nicht solide.}
\end{minipage}
\end{center}
\end{forschung}

\begin{merksatz}
\boxtitel{Was Sie jetzt können}
\begin{center}
\renewcommand{\arraystretch}{1.5}
\begin{tabular}{@{}P{3.0cm}P{5.6cm}@{}}
\toprule
\textbf{Condorcet} & Mehrheiten sind nicht übertragbar\\
\textbf{Die Tagesordnung} & entscheidet vollständig\\
\textbf{Satz 12.1} & Arrow: es gibt kein Verfahren\\
\textbf{Satz 12.2} & keines ist manipulationsfest\\
\textbf{Satz 12.3} & eingipflig: der Median gewinnt\\
\bottomrule
\end{tabular}
\end{center}

\begin{center}
\begin{minipage}{0.88\textwidth}
\kursiv{Arrow verlangt eine vollständige Reihung als Ergebnis. Das nächste Kapitel verlangt weniger --- nur eine Zuteilung --- und findet Verfahren, die tatsächlich funktionieren.}
\end{minipage}
\end{center}
\end{merksatz}
